% ----------------------------------------------------------------------
%
%%                          Principal.tex
%
%----------------------------------------------------------------------
%
%  El archivo "Principal.tex" es el documento principal o maestro que
%  reúne la configuración del entorno LaTeX, el formato  del documento,
%  paquetes de configuración, datos de alumnnos y jurado, así como
%  incluir los archivos que contienen la información de cada capítulo
%  o sección.
%
%----------------------------------------------------------------------
%
% Los archivos utilizados para este documento son:
%
%       preliminares/*.tex : documentos preliminares del documento
%       dedicatoria.tex, agradecimientos.tex, resumen.tex, abstract.tex,
%       introducción.tex, nomenclatura.tex
%       capitulos/*.tex : capítulos del documento
%       marcoreferencia.tex, diseconceptual.tex, disedetallado.tex
%       implementacón.tex, discusion.tex, conclusiones.tex
%       apendices/*.tex: apéndices y anexos del documento
%       Principal.tex: Archivo maestro del documento desde el cual se realiza
%       la compilación
%       upiita.cls : Archivo de la clase de documento, contiene las definiciones
%       y estructura del formato
%       upiitatesis.sty Configuración y definiciones del documento
%       XBiblioteca.bib: Archivo de las referencias
% ---------------------------------------------------------------------
%       Inicio del documento se define la clase y
%           UpiiTeXis: Formatos de Reporte de Trabajo termina o Proyecto Terninal UPIITA-IPN
%   Copyright (C) 2023 Juan Carlos Guzmán Salgado
%                      Griselda Sánchez Otero
%                      Diego Alonso Flores Hernández
%%%%%%%%%%%%%%%%%%%%%%%%%%%%%%%%%%%%%%%%%%%%%%%%%%%%%%%%%%%%%%%%%%%%%%%%%%%%
% Este trabajo está desarrollado con base en la idea de  TeXiS pero no utliza
% sus funciones, estilos o clases, solo la estructura de archivos
% Copyright 2009 Marco Antonio Gomez-Martin, Pedro Pablo Gomez-Martin
%
% The complete last TeXiS package can
% be obtained from http://gaia.fdi.ucm.es/projects/texis/
%%%%%%%%%%%%%%%%%%%%%%%%%%%%%%%%%%%%%%%%%%%%%%%%%%%%%%%%%%%%%%%%%%%
%  The UpiiTeXis template itself is distributed under the
% conditions of the LaTeX Project Public License
% (http://www.latex-project.org/lppl.txt), the manual content
% uses the CC-BY-SA license that stays that you are free:
%
%    - to share & to copy, distribute and transmit the work
%    - to remix and to adapt the work
%
% under the following conditions:
%
%    - Attribution: you must attribute the work in the manner
%      specified by the author or licensor (but not in any way that
%      suggests that they endorse you or your use of the work).
%    - Share Alike: if you alter, transform, or build upon this
%      work, you may distribute the resulting work only under the
%      same, similar or a compatible license.
%
% The complete license is available in
% http://creativecommons.org/licenses/by-sa/3.0/legalcode
%
%---------------------------------------------------------------------
\typeout{Copyright Juan Carlos Guzmán Salgado Griselda Sánchez Otero   Diego Alonso Flores Hernández}
\documentclass[letterpaper,11pt]{upiita}
\pagenumbering{arabic}  % P08: define \thepage antes de hyperref, para que el visor muestre i, ii... y 1, 2...
\usepackage{configuracion/upiitatesis}
\renewcommand{\rmdefault}{lmr}   % texto en versión vectorial de la misma fuente (P29: sin cm-super salía en mapa de bits)
\renewcommand{\ttdefault}{lmtt}  % ídem para \verb y texto monoespaciado
% P15: en portada y acta, \bf + \fontseries{sb} pide Montserrat "sbx", que no existe; se declara
% como la negrita que LaTeX ya usaba en su lugar (mismo aspecto, sin avisos)
{\fontencoding{T1}\fontfamily{Montserrat-TLF}\selectfont}% carga las definiciones de la fuente
\DeclareFontShape{T1}{Montserrat-TLF}{sbx}{n}{<->ssub * Montserrat-TLF/bold/n}{}
\usepackage{../revision}  % marcas de revisión (Escrito/CODIGOS_REVISION.md); para entregar: [final]
%%%%%%%%%%%%%  Datos de la Portada y acta
\title{Sistema de localización acústica multimicrófono basado en mecanismos de detección vibracional de escorpiones para la mejora de la inteligibilidad del habla en entornos acústicamente adversos}
\unidad{\Upiita}
\materia{Trabajo Terminal I}  % ya no se imprime: la portada y el acta muestran "UPIITA" (P31)
\academia{Bi\'onica}
\grado{Ingeniero en Bi\'onica}      %%%% grado
\mes{Ciudad de México, a 21 de enero de}                               %%%% mes
\anio{2027}                              %%%% año
%\impreso{pdf}							%%%% opciones: pdf, papel
%%%%%%%%%%%%  Nombre del Presidente y Secretario del jurado
\Presidente{Dr. Rafael Bayareh Mancilla}
\Titular{Dra. Lilia Martínez Pérez}
%%%%%%%%%%%%%%%%%%%%%%% Datos del autor%%%%%%%%%%%%%%%%%%%%%%%%%%%%%%
%%  para varios asesores (máximo 4) y autores (alumnos)
%%  dejar el campo en blanco o borrar los que no deban aparecer
%%  Cuando el número de asesores=1 sólo se toma en cuenta la variable
%%% alumnoa,
%%  Cuando el número de asesores=2 se toman alumnoa y asesorb, etc
%%%%%%%%%%%%%%%%%%%%%%%%%%%%%
\nalumnos{2}
%%%%%%% dependiendo del numero de alumnos máximo 4,
%%%%%%%  eliminar o comentar la líneas, 1 alumno  =alumnoa
\alumnoa{Deloiza Rodríguez Arturo}
\alumnob{Jiménez Jiménez Atenea}
\alumnoc{Jan Kaarle iohannes Karl}
\alumnod{Giovanni Karol iohannes Karl}
%%%%%%%%%% número de asesores 1= solo se toma en cuenta la variable asesora,
%%%%%%%%%%  2= se toman asesora y asesorb
\nasesores{3}
%%%%%%% en caso de un único asesor asesora,
%%%%%%%  eliminar o comentar las líneas
\asesora{Ing. Carlos Ríos Ramírez}
\asesorb{Dr. Álvaro Anzueto Ríos}
\asesorc{Dr. Carlos Carrizales Velazquez}
%%%%%% Inicio de documento %%%%%%
\begin{document}
\makeatletter \def\@roman#1{\romannumeral #1}\makeatother
\frontmatter
%%%%% Inicio de Sección de preliminares
%%Preliminar
%%creacion de carátula y hoja de actas
\portada                  %% Portada s
\acta                     %% Portada con firmas
\dedicatoria              %%%modificar archivo dedica.tex
\agradecimientos          %%%modificar archivo agradece.tex
\tabladecontenido
\chapter{Resumen}


\textbf{Resumen:}



\textbf{Palabras Clave:} Robot, Bípedo, Humanoide,  Marcha humana, Criterio de equilibrio, Punto de momento Cero, Actuadores, Sensores, Fuerza, Par torsor, Modelo matemático, Estructura, Diseño mecánico, Esfuerzos, Desplazamientos, Optimización, Simulaciones, Trayectoria, Cinemática directa e inversa, Análisis. \\



%%%%%%fin del archivo

\endinput                %%%modificar archivo preliminares/resumen.tex
\chapter{Abstract}

\textbf{Abstract:}




\textbf{Palabras Clave:} Robot, Bípedo, Humanoide,  Marcha humana, Criterio de equilibrio, Punto de momento Cero, Actuadores, Sensores, Fuerza, Par torsor, Modelo matemático, Estructura, Diseño mecánico, Esfuerzos, Desplazamientos, Optimización, Simulaciones, Trayectoria, Cinemática directa e inversa, Análisis. \\


%%%%%%fin del archivo
\endinput                %%%modificar archivo preliminares/abstract.tex
\listadefiguras
\listadetablas
%%%%%%%%%%%%%%%%%%%%%%%%%%%%%%%%%%%%%%%%%%%%%%%%%%%%%%%%%%%%%%%%%%%%%%%%%%%%%%%%%%%%%%%%%%%%%%%
\nomenclatura

%\printnomenclature





%%%%%%%%  LISTA DE ACRÓNIMOS  % % A = GRUPO Acrónimos  %%%%%
%%	\nomenclature[A]{Acrónimo}{Definición}
\nomenclature[A]{SEMARNAT}{Secretar\'ia de Medio Ambiente y Recursos Naturales}
\nomenclature[A]{RSU}{Residuos Solidos Urbanos}
\nomenclature[A]{NOM}{Norma Oficial Mexicana}
\nomenclature[A]{FBS}{Estructura Funcional de Desglose\textit{(por sus siglas en ingl\'es Functional Breakdown Structure )}}
\nomenclature[A]{IDEF-0}{Definici\'on de integraci\'on por modelado de funciones\textit{(por sus siglas en ingl\'es Functional Breakdown Structure )}}
\nomenclature[A]{SPA }{Arquitectura F\'isica del sistema\textit{(por sus siglas en ingl\'es System Physical Architecture )}}
\nomenclature[A]{AHP}{Proceso anal\'itico jer\'arquico \textit{(del ingl\'es Analytic Hierachy Process)}}


%%%%%%%%% SIMBOLOS % B = GRUPO Símbolos  %%%%%
%\nomenclature[B]{símbolo}{Definición}

\nomenclature[B]{$\rho_{2}$}{Relaci\'on de transmisi\'on entre la catarina conducida y la catarina motriz}
\nomenclature[B]{$N_{1}$}{N\'umero de dientes de la catarina motriz}
\nomenclature[B]{$N_{2}$}{N\'umero de dientes de la catarina conducida}
\nomenclature[B]{$K_{t}$}{factor de concentraci\'on de esfuerzos geom\'etricos}
\nomenclature[B]{$K_{f}$}{Factor de correci\'on por concentraci\'on de esfuerzos}
\nomenclature[B]{$F_{d}$ }{actor del di\'ametro del transportador}
\nomenclature[B]{$f_{b}$ }{Factor del buje para colgante}
\nomenclature[B]{$F_{f}$}{Factor de helicoidal}
\nomenclature[B]{$F_{m}$}{Factor de material}
\nomenclature[B]{$F_{p}$ }{Factor de paletas}
\nomenclature[B]{$F_{fo}$}{Factor de sobrecarga}
\nomenclature[B]{$e$}{Eficiencia de la transmisi\'on}
\nomenclature[B]{$a$}{The number of angels per unit area}%
\nomenclature[B]{$N$}{The number of angels per needle point}%
\nomenclature[B]{$A$}{The area of the needle point}%
\nomenclature[B]{$\sigma$}{The total mass of angels per unit area}%
\nomenclature[B]{$m$}{The mass of one angel}
\nomenclature[B]{$\beta$}{The second letter of the greek alphabet}
\nomenclature[B]{$\alpha$}{The first letter of the greek alphabet}
\nomenclature[B]{$\mathcal {L}$}{ Transformada de Laplace}
\nomenclature[B]{$\mathcal {F}$}{ Transformada de Fourier}

%%%%% VARIABLERS FÍSICAS % C = GRUPO Variables físicas  %%%%%
%   Las variables físicas contienen unidades
%\nomenclature[C]{variable)}{{descripción}\nomunit{\SI[group-digits=false]{valor}{unidades}}}

\nomenclature[C]{$F_c$}{{Fuerza de corte}\nomunit{\SI[group-digits=false]{}{\N}}}
\nomenclature[C]{$\tau_{c}$}{{ Esfuerzo cortante}\nomunit{\SI[group-digits=false]{}{\MPa}}}
\nomenclature[C]{$\tau_{v}$}{{Esfuerzo variable}\nomunit{\SI[group-digits=false]{}{\MPa}}}
\nomenclature[C]{$\tau_{m}$ }{{ Esfuerzo medio}\nomunit{\SI[group-digits=false]{}{\MPa}}}
\nomenclature[C]{$e$}{{Espesor de l\'amina}\nomunit{\SI[group-digits=false]{}{\m}}}
\nomenclature[C]{$L$}{{Longitud del borde de corte}\nomunit{\SI[group-digits=false]{}{\m}}}
\nomenclature[C]{$M_{t}$}{{Momento de torsi\'on del eje de trituraci\'on }\nomunit{\SI[group-digits=false]{}{\N\m}}}
\nomenclature[C]{$Mt_{v}$}{{Momento de torsi\'on variable }\nomunit{\SI[group-digits=false]{}{\N\m}}}
\nomenclature[C]{$Mt_{m}$}{{Momento de torsi\'on medio}\nomunit{\SI[group-digits=false]{}{\N\m}}}
\nomenclature[C]{$\omega$}{{Velocidad angular}\nomunit{\SI{}{\radian\per\second}}}
%%%%%%fin del archivo
\endinput           %%%modificar archivo preliminares/nomenclatura.tex
%% En esta sección se agregan Símbolos
%%
%%  \nomenclature[Grupo]{acrónimo}{definición}
%%            A = GRUPO Acrónimos
%%            B = GRUPO Símbolos
%%            C = GRUPO Variables físicas
%%%%%%%%%%%%%%%%%%%%%%%%%%%%%%%%%%%%%%%%%%%%%%%%%%%%%%%%%%%%%%%%%%%%%%%%%%%%%%%%%%%%%%%%%%%%%%%%
\Introduccion                                 %%%modificar archivo preliminares/introduccion.tex
%%%% Contiene los siguientes elementos
%% Definición del problema
%% Justificación
%% Objetivo
%%      %%General
%%      %%Específico
%% enfoque mecatrónico
%% Antecedentes
%% Capitulado descripción general del documento mencionando contenidos de cada
%%              capitulo ;Descripción de los capítulos
%%%%%%%%%%%%%%%%%%%%%%%%%%%%%%%%%%%%%%%%%
%%%%%% Inicia Capítulos del trabajo%%%
%%%%%%%%%%%%%%%%%%%%%%%%%%  capitulo 1 %%%%%%%%%%%%%%%%%%%%%%%%%%%
\MarcodeReferencia
%%% modificar archivo capitulos/marcoreferencia.tex
%%%%%%%%%%%%%%%%%%%%%%%%%%capitulo1111111111%%%%%%%%%%%%%%%%\media
%\media  %%ajuste de contador de TOC %no eliminar
%\chapter{Marco de referencíal}

\section{Marco teórico}


\subsection{Metodología mecatrónica}

\subsection{ \dots  \dots ...}


\section{Marco procedimental}
%\newgeometry{margin=1cm,includefoot}
%\layout
%\printinunitsof{cm}
%\pagevalues
\subsection{ \dots  \dots ...}


%%%%%%fin del archivo
\endinput 
%%%modificar archivo capitulos/marcoreferencia.tex
%%%%%%%%%%%%%%%%%%%%%%%%%%%%%%%%%%%%%%%%%%%%%%%%%%%%%%%%%%%%%%%%%%%
%%%%%%%%%%%%%%%%%%%%%%%%%%%%%%%%%%%%%%%%%%%%%%%%%%%%%%%%%%%%%%%%%%%
% !TeX encoding = UTF-8
\chapter{Diseño del sistema}
%
\section{Diseño Conceptual}
%
\subsection{Necesidades - Requerimientos}
%
\subsection{Arquitectura funcional (funciones)}

\subsection{Arquitectura física (módulos)}

\subsection{Propuestas de solución}

\subsubsection{Módulos (M$\textsubscript{1}$  \dots M$\textsubscript{n}$)}

\subsubsection{Integración (sistema)}


\subsection{Validación}

\subsection{Selección diseño conceptual}



%%%%%%fin del archivo
\endinput
%%%modificar archivo capitulos/diseconceptual.tex
%%%%%%%%%%%%%%%%%%%%%%%%%%%%%%%%%%%%%%%%%%%%%%%%%%%%%%%%%%%%%%%%%%%
%%%%%%%%%%%%%%%%%%%%%%%%%%%%%%%%%%%%%%%%%%%%%%%%%%%%%%%%%%%%%%%%%%%
% !TeX encoding = UTF-8
\section{Diseño Detallado}

\subsection{Detalle módulo 1 (M{l})}

\subsubsection{Validación (M{l})}

\subsection{Detalle módulo 2 (M{2})}

\subsubsection{Validación (M{2})}


\subsection{Detalle módulo n (M{n})}

\subsubsection{Validación (M{n})}

% [DUDA] P26: títulos de la plantilla para Mecatrónica; el protocolo es de Biónica (enfoque biomimético)
\subsection{Integración sistema mecatrónico}
\duda{Título de la plantilla para Mecatrónica: no corresponde al protocolo de Biónica. Adaptar o quitar.}

\subsection{Validación sistema mecatrónico}
\duda{Título de la plantilla para Mecatrónica: no corresponde al protocolo de Biónica. Adaptar o quitar.}



%%%%%%fin del archivo
\endinput 
%%%modificar archivo capitulos/disedetallado.tex
%%%%%%%%%%%%%%%%%%%%%%%%%%%%%%%%%%%%%%%%%%%%%%%%%%%%%%%%%%%%%%%%%%%
\chapter{Implementación del sistema}
%
\section{Implementación módulo 1 (M{l})}

\subsubsection{Verificación (M{l})}


\section{Implementación módulo 2 (M{2})}

\subsubsection{Verificación (M{2})}


.  \dots ...
\section{Implementación módulo 2 (M{n})}

\subsubsection{Verificación (M{n}}


\section{Implementación módulo 1 (M{l})}

\subsubsection{Verificación (M{l})}





%%%%%%fin del archivo
\endinput 
%%%modificar archivo capitulos/implementacion.tex
%%%%%%%%%%%%%%%%%%%%%%%%%%%%%%%%%%%%%%%%%%%%%%%%%%%%%%%%%%%%%%%%%%%
%%  \final
\chapter{Discusión / Análisis de resultados}

\subsection{Análisis de ingeniería}


\subsection{Análisis de costos}


\subsection{Análisis de valor}



%%%%%%fin del archivo
\endinput
%%%modificar archivo capitulos/discusion.tex
%%%%%%%%%%%%%%%%%%%%%%%%%%%%%%%%%%%%%%%%%%%%%%%%%%%%%%%%%%%%%%%%%%%
%%%Ajusta el contador de secciones y subsecciones para ete capítulo
\setcounter{secnumdepth}{-1}
%%%%%%%%%%%%%%%%%%%%%%%%%%%%%%%%%%%%%%%%%%%%%%%%%%%%%%%%%%%%%%%%%%%

\chapter{Conclusiones}

una conclusi\'on\index{conclusión}
\cite{AR17}

\section{Recomendaciones y trabajo futuro}

%



%%%%%%fin del archivo
\endinput 
%%%modificar archivo capitulos/conclusiones.tex
%%%%%%%%%%%%%%%%%%%%%%%%%%%%%%%%%%%%%%%%%%%%%%%%%%%%%%%%%%%%%%%%%%%
\newpage
%%%%%%%%%%%%%%%%%%%%%%%%%%%%%%%%%%%%%%%%%%%%%%%%%%%%%%%%%%%%%%%%%%
%%%%%%%%%%%%%%%%%%%%%%%  BIBLIOGRAFIA  %%%%%%%%%%%%%%%%%%%%%%%%%%%
%%%   \adicional
\referencias
%%%%%%%%%%%%%%%%%%%%%%%%%%%%%%%%%%%%%%%%%%%%%%%%%%%%%%%%%%%%%%%%%%%
%%%%%%%%%%%%%%%%%%%%%%%%%    APENDICES  %%%%%%%%%%%%%%%%%%%%%%%%%%%
\backmatter
\part{Apéndices}
\appendix
\chapter{Ap\'endice 1}\label{apendice1}
\stepcounter{chapter} %instrucción necesaria para cambiar numeración

Incluyen información que ayuda a interpretar parte del contenido del libro, o aspectos más técnicos y menos esenciales del libro, como información complementaria.


\bigskip

Ab hier beginnt der \verb|Anhang|.






%%%%%%fin del archivo
\endinput 
\chapter{Der Zweite Anhang (Ap\'endice 2)}\label{apendice2}
\stepcounter{chapter} %instrucción necesaria para cambiar numeración



Incluyen información que ayuda a interpretar parte del contenido del libro, o aspectos más técnicos y menos esenciales del libro, como información complementaria.


\bigskip





%%%%%%fin del archivo
\endinput 
\chapter{Der dritte Anhang (Ap\'endice 3)}\label{apendice3}
\stepcounter{chapter} %instrucción necesaria para cambiar numeración


Incluyen información que ayuda a interpretar parte del contenido del libro, o aspectos más técnicos y menos esenciales del libro, como información complementaria.

\bigskip




%%%%%%fin del archivo
\endinput 
\part{Anexos}
\anexo
\chapter{Anexo 1. Hoja de datos}\label{anexo1}
\stepcounter{chapter} %instrucción necesaria para cambiar numeración

Una ficha técnica, hoja técnica u hoja de datos (datasheet en inglés), también ficha de características u hoja de características, es un documento que resume el funcionamiento y otras características de un componente (por ejemplo, un componente electroóico) o subsistema (por ejemplo, una fuente de alimentación) con el suficiente detalle para ser utilizado por un ingeniero de diseño y diseñar el componente en un sistema.

Comienza típicamente con una página introductoria que describe el resto del documento, seguido por los listados de componentes específicos, con la información adicional sobre la conectividad de los dispositivos. En caso de que haya código fuente relevante a incluir, se une cerca del extremo del documento o se separa generalmente en otro archivo.

Las fichas técnicas no se limitan solo a componentes electrónicos, si no que también se dan en otros campos de la ciencia, como por ejemplo compuestos químicos o alimentos.

\bigskip






%%%%%%fin del archivo
\endinput 
\chapter{Anexo 2. Hoja de datos}\label{anexo2}
\stepcounter{chapter} %instrucción necesaria para cambiar numeración

Una ficha técnica, hoja técnica u hoja de datos (datasheet en inglés), también ficha de características u hoja de características, es un documento que resume el funcionamiento y otras características de un componente (por ejemplo, un componente electroóico) o subsistema (por ejemplo, una fuente de alimentación) con el suficiente detalle para ser utilizado por un ingeniero de diseño y diseñar el componente en un sistema.

Comienza típicamente con una página introductoria que describe el resto del documento, seguido por los listados de componentes específicos, con la información adicional sobre la conectividad de los dispositivos. En caso de que haya código fuente relevante a incluir, se une cerca del extremo del documento o se separa generalmente en otro archivo.

Las fichas técnicas no se limitan solo a componentes electrónicos, si no que también se dan en otros campos de la ciencia, como por ejemplo compuestos químicos o alimentos.




\bigskip



%%%%%%fin del archivo
\endinput 
\chapter{Anexo 3 (Anexo título 3)}\label{anexo3}
\stepcounter{chapter} %instrucción necesaria para cambiar numeración

Una ficha técnica, hoja técnica u hoja de datos (datasheet en inglés), también ficha de características u hoja de características, es un documento que resume el funcionamiento y otras características de un componente (por ejemplo, un componente electrónico) o subsistema (por ejemplo, una fuente de alimentación) con el suficiente detalle para ser utilizado por un ingeniero de diseño y diseñar el componente en un sistema.

Comienza típicamente con una página introductoria que describe el resto del documento, seguido por los listados de componentes específicos, con la información adicional sobre la conectividad de los dispositivos. En caso de que haya código fuente relevante a incluir, se une cerca del extremo del documento o se separa generalmente en otro archivo.

Las fichas técnicas no se limitan solo a componentes electrónicos, si no que también se dan en otros campos de la ciencia\index{ciencia}, como por ejemplo compuestos químicos o alimentos.

\bigskip





%%%%%%fin del archivo
\endinput 
\newpage
%%%%%%%%%%%%%%%%%%%%%%%%%%%%%%%%%%%%%%%%%%%%%%%%%%%%%%%%%%%%%%%%%%%%%%%
%%%%%%%%%%%%%%%%%%%%%%%%%%%%  INDICE ALFABÉTICO%%%%%%%%%%
\IndiceAlfabetico
\end{document}
%%%%%%%%%%%%%%%%%%%%%%%%%%%%%%%%%%%%%%%%%%%%%%%%%%%%%%%%%%%%%%%%%%%%%%%%%%%%%%%%%% 